\section{orientifold の refined 指数と全 $d$ の構造}\label{sec:index}

\subsection{設定}
AgSc \cite{AgSc12} §3 にしたがい、M 理論を $X \times \mathrm{TN} \times S^1$ 上で $(q, t)$ で捻る。orientifold の作用 $\tilde I$ は、$X$ の反正則な対合 $I$ と、時空の $z_2 \to -z_2$ の組である。
結び目 $K$ の probe は M5 brane で、$z_2$ 平面を張る（t-brane。開弦の分母が $1/(1-t)$ であることによる）。被覆には $L_K$ と像 $I(L_K)$ の 2 枚がある。
被覆の一粒子状態の空間は $H_1 = H_{\mathrm{int}} \otimes F_{z_2}$ である。$F_{z_2}$ は $z_2$ 方向のモード $t^l$（$l \ge 0$）から成る。$\tilde I$ は $F_{z_2}$ に $(-1)^l$ で働く。

\begin{hypothesis}[一粒子状態の分類]\label{hyp:classes}
商の一粒子状態は被覆の $\tilde I$ の軌道で、AgSc の積表示 (3.13) のとおり次の 3 種に分かれる。
\begin{enumerate}
\item $\tilde I$ で入れ替わる対 $\{\psi, \tilde I\psi\}$：全モードの塔をもつ。
\item $\tilde I\psi = \varepsilon\psi$（$\varepsilon = \pm 1$）：$(-1)^l\varepsilon = 1$ のモードだけをもつ。
\item $z_2 = 0$ に凍結した状態：モードの塔をもたない。
\end{enumerate}
\end{hypothesis}

\subsection{全 $d$ の構造}
% [O1][O2]
\begin{proposition}\label{prop:alld}
仮説 \ref{hyp:classes} のもとで、状態の電荷を $m$ とすると
\[
g = f - \tfrac12 f^{\cov} = \frac12\,\frac{E}{1+t} + C,\qquad E = \sum_{\text{不変}} \varepsilon_\psi\, m_\psi,\qquad C = \sum_{\text{凍結}} m_\psi .
\]
\end{proposition}
\begin{proof}
被覆の指数は $f^{\cov} = \sum_{\text{対}} 2m/(1-t) + \sum_{\text{不変}} m/(1-t)$ である。
商の指数は $f = \sum_{\text{対}} m/(1-t) + \sum_{\text{不変}} m \sum_{(-1)^l\varepsilon = 1} t^l + \sum_{\text{凍結}} m$ である。
$\sum_{(-1)^l\varepsilon = 1} t^l = \tfrac12[1/(1-t) + \varepsilon/(1+t)]$ を代入すれば主張が従う。
\end{proof}

\begin{remark}\label{rem:Esector}
$\tilde I$ で不変な状態は、$\tilde I$ で自分自身に移るセクター（brane と image を同数もつセクター）にしか無い。したがって $E$ はそのセクターにだけ現れる（命題 \ref{prop:boundary} の系も参照）。
$\tilde I$ が brane と image の入れ替えとして働くなら、そのトレースは冪和の基底で $p_2$ 型の方向（Adams の $\psi^2$）に入り、$E_{(2)} = -E_{(1,1)}$ となる。
後半は Hopf 絡み目（計算結果 \ref{cr:E22}）と結び目の 2 箱（計算結果 \ref{cr:largeN}）で確かめた性質である。
\end{remark}

\subsection{trace 表示：Witten の二つの公式の SO$(2N)$ 版}
Witten \cite{Witten11} は、Chern--Simons の分配関数を 5 次元（6 次元 $(2,0)$ 理論）の $Q$ コホモロジー上のトレース $Z^{\mathrm{CS}} = \mathrm{Tr}_{\mathcal K}\, q^{P}(-1)^{F}$ と書き（§4.1）、単純レースでない群については、大きい理論の空間で外部自己同型 $\zeta$ を挿入したトレースと、固定された理論の普通のトレースという二つの表示を与えた（§3.4）。
D 型の $(2,0)$ 理論は、M5 brane を $X/\mathbb{Z}_2$（$\mathbb{Z}_2$ は法束に $-1$）に巻いて得られる（同論文 §5.1.4、\cite{AgSc12} 付録 A）。これが本論文の $\tilde I$ の 5-brane 側の起源である。
% [O9]
\begin{proposition}\label{prop:trace}
仮説 \ref{hyp:classes} のもとで、商の一粒子の指数は被覆の射影 $P_+ = (1 + \tilde I)/2$ と凍結した状態の和で書け、
\[
g = f - \tfrac12 f^{\cov} = \tfrac12\,\mathrm{Tr}_{H_1^{\cov}}\bigl(\tilde I\cdot(\text{指数の重み})\bigr) + C
\]
である。右辺の第 1 項は $\tfrac12 E/(1+t)$ に等しい。
\end{proposition}
\begin{proof}
$f = \mathrm{Tr}_{H_1^{\cov}}(P_+\cdot) + \sum_{\text{凍結}} m$ と、$\tilde I$ が不変な状態に $\varepsilon$、$z_2$ のモードに $(-1)^l$ で働くことから、$\mathrm{Tr}(\tilde I\cdot) = \sum_{\text{不変}}\varepsilon m\sum_l (-t)^l = E/(1+t)$。
\end{proof}
$E$ は Witten の「$\zeta$ を挿入したトレース」の型、$C$ は商の理論にしかない状態の寄与である。命題 \ref{prop:EC} の不定性は、一つの数を二つの表示に分ける仕方の不定性と読める。

\subsection{境界ゼロモードへの $\tilde I$ の作用}
KN \cite{KN17} §3.2 では、境界 $h$ 個の曲線の境界ゼロモード $d\theta_i$（$\sum_i d\theta_i = 0$）が hook 表現を張り、t-brane 枠では 1 モードあたり $(1 - q^{-1})$（規格化 $q^{1/2}$）を与える。
% [O10]
\begin{proposition}\label{prop:boundary}
KN の境界ゼロモードの描像を仮定し、$\tilde I$ が $L$ の境界と $I(L)$ の境界を入れ替えるとする。$\tilde I$ 不変な曲線は $L$ と $I(L)$ の境界を同数 $n$ 組もち、対ごとに和 $d\theta_i + d\theta_i'$ は偶、差 $d\theta_i - d\theta_i'$ は奇である。拘束 $\sum d\theta = 0$ は偶のモードを一つ減らす。したがって
\[
E\ \text{の境界因子} = z_q^{\,n-1}(z_q^+)^{\,n},\qquad
P := (1-t)f^{\cov}_{\text{固定セクター}}\ \text{の境界因子} = z_q^{\,2n-1}.
\]
系として、$E$ は brane と image の境界を同数もつセクター（$p_2$ 型の方向）にだけ現れる。
\end{proposition}
\begin{proof}
奇のモードでは、$\tilde I$ を挿入したトレースで $(1 - q^{-1})$ が $(1 + q^{-1})$ に替わる。偶 $n-1$ 個、奇 $n$ 個を掛け合わせればよい。
\end{proof}
% [O10]
\begin{compresult}\label{cr:boundary}
保存済みのデータで、命題 \ref{prop:boundary} の予言がすべて一致した。
\begin{center}
\begin{tabular}{llll}
\toprule
対象 & $n$ & $E$ の境界因子 & $P$ の境界因子 \\
\midrule
Hopf $(2,2)$, $(2,4)$ & 2 & $z_q(z_q^+)^2$ & $z_q^3$ \\
$T(2,3)$, $T(2,5)$ の $p_2$ 方向 & 1 & $z_q^+$（整数計画は $z_q^+$ でなければ解なし） & $z_q$ \\
framed unknot（framing 2, 4） & 1 & $z_q^+$ & $z_q$ \\
\bottomrule
\end{tabular}
\end{center}
\end{compresult}
計算結果 \ref{cr:E22} の $E_{22}$ の $(z_q^+)^2$ と、§\ref{sec:positivity} の「結び目の境界因子は和でなければならない」は、この命題の帰結である。残る未知は、モジュライ空間の部分 $H^*(\mathcal{M}_{g,h,\beta})$ への $\tilde I$ の作用（$\beta$ ごとの符号 $\varepsilon$）である。

\subsection{$\sigma$ 対称性の起源}
$\sigma\colon q\mapsto t^{-1}, t\mapsto q^{-1}$ は、KN の変数 $q_\ell = (qt)^{1/2}$, $q_r = (q/t)^{1/2}$ では $q_\ell \mapsto q_\ell^{-1}$、$q_r$ 固定、すなわち Jacobian の Lefschetz 作用 SU$(2)_\ell$ の Weyl 反射である。
KN の $\hat f$ の $\sigma$ 不変性は、取っ手のゼロモードが SU$(2)_\ell$ の二重項をなすこと（取っ手あたり $-(q_\ell + q_\ell^{-1}) + (q_r + q_r^{-1}) = G$）から来る。
凍結した状態は向き付けられない面 $N = \hat\Sigma/I$ に巻いた M2 で、$N$ には複素構造が無く $H^1(N)$ に Lefschetz の $\mathfrak{sl}_2$ は働かない。したがって $C$ に $\sigma$ 対称性を期待する理由は無く、多箱の $c = 1$ 部分が $G$ 展開できないこと（§\ref{sec:generalN}）はこれで説明できる（予想）。[1] の $c = 1$ 部分が $\sigma$ 対称であることは、次の観察で説明できる（予想）。[1] の $c = 1$ 部分の $\nu^\beta$ のブロックの $G$ の次数は、brane--image のセクター $(t^{-1/2}-t^{1/2})x/z_q$ の $\nu^{2\beta}$ のブロックの $G$ の次数のちょうど半分である（$T(2,3)$, $T(2,5)$ のすべての $\beta$ で確認）。[1] の凍結した状態は brane--image の環状の弦の商（Möbius の帯）で、商の取っ手 $g$ 個が被覆の取っ手 $2g$ 個に当たり、商の取っ手が $G$ を担う。多箱の $c = 1$ 部分が $\sigma$ 対称でない理由は未解明である。

Mariño の $\sum_{d\ \mathrm{odd}}$ の定義 $g^{\mathrm{odd}}$ とは、$g_{ab} = g^{\mathrm{odd}}_{ab} - \tfrac12 g^{\mathrm{odd}}_{a/2, b/2}(x^2)$ で結ばれる。

% [O4]
\begin{compresult}
Hopf 絡み目（$a, b \le 4$、§\ref{sec:hopf}）で、全 $d$ の $g_{ab}$ の分母は（偶, 偶）のセクターで $2(1+t)\times$ 単項式、それ以外で単項式である。
$\sum_{d\ \mathrm{odd}}$ の定義で現れる $4(1+t)(1+t^2)$ は、全 $d$ の定義では打ち消す。
\end{compresult}

\subsection{$E$ と $C$ の分離の不定性}
% [O3]
\begin{proposition}\label{prop:EC}
整数係数の Laurent 多項式 $X$ について、$E \to E + 2(1+t)X$, $C \to C - X$ としても $g$ は変わらない。したがって $g$ から決まるのは $E$ の $2(1+t)$ を法とした類だけで、特に $E$ そのものは決まらない（$E(-1) = [2(1+t)g]_{t=-1}$ は決まる）。
\end{proposition}
\begin{proof}
$\tfrac12 \cdot 2(1+t)X/(1+t) = X$ なので、$E$ の変化は $C$ の変化で打ち消す。
状態の言葉では、次の二つの状態の組は、商の指数 $f$ と被覆の指数 $f^{\cov}$ の両方で区別できない。
\begin{itemize}
\item 凍結した状態 $m$ と、$\tilde I$ で入れ替わる対 $\{m\}$, $\{mt\}$。
\item $\varepsilon = +1$ の不変な状態 $m$, $m$, $mt$, $mt$。
\end{itemize}
実際、$f^{\cov}$ への寄与はどちらも $2(m + mt)/(1-t)$、$f$ への寄与はどちらも $2m/(1-t)$ である。
\end{proof}

% [O7]
\begin{corollary}\label{cor:Efree}
冪和の $p_1^n$ 方向と、箱の数が奇数のセクター（$E = 0$）は、$E$ によらない。これらの方向では、正値性を $E$ の不定性によらずに判定できる。
\end{corollary}

\subsection{閉弦での確認}
% [O5]
\begin{compresult}\label{cr:closed}
D$_N$ の積表示 $Z = \prod_{m<\beta}\prod_{\alpha>0}(q^{-m/2}t^{-(\alpha,\rho)/2} - q^{m/2}t^{(\alpha,\rho)/2})$（$t = q^\beta$）から作った $-\log Z$ の $d$ 番目と、AgSc §6 の 3 項
\[
\frac{1}{2d}\frac{Q^d t^{d/2} q^{-d/2}}{z_q(q^d) z_t(t^d)} + \frac1d \frac{(Q t/q)^{d/2}}{z_q(q^d)} - \frac{1}{2d}\frac{Q^d t^{d/2}q^{-d/2}}{z_q(q^d)(t^{d/2} + t^{-d/2})},\qquad Q = t^{2N-1}
\]
との差は、$N$ の 1 次式（古典的な項と定数写像）だけである（$N = 2, \dots, 7$）。
\end{compresult}
AgSc の 3 項は、命題 \ref{prop:alld} の閉弦版である。P$^1$ は $\varepsilon = +1$ の不変な状態で、第 1 項と第 3 項の和が $z_2$ の偶数のモードだけの寄与になる。RP$^2$ は凍結した状態で、第 2 項に当たる。

Avetisyan--Mkrtchyan \cite{AM22} は、同じ分配関数を多重サイン関数で書き、D・C 型の自由エネルギーを 3 種に整理した。奇数の $n$ の $c = 1$ の項、全 $n$ の「Klein 瓶」（$c = 2$）の項、Kähler 変数の「三重化」である。
% [O6]
\begin{proposition}\label{prop:AM}
AM の変数 $(Q, T) = (q^{1/2}, t^{1/2})$ で、$X = e^{\tau_1/2}$, $e^{\tau_2} = X^2/Q$, $e^{\tau_0} = X^2 T/Q$ とする。このとき $n$ ごとの有理式として、次の恒等式が成り立つ。
\[
\underbrace{-\frac{X^{2n}}{2n(Q^{2n} - Q^{-2n})}}_{\text{奇数 $n$ に直す差}}
+ \underbrace{\frac{e^{n\tau_2}}{2n}\frac{Q^n T^{-n} - T^n Q^{-n}}{(Q^{2n} - Q^{-2n})(T^n + T^{-n})}}_{\text{Klein 瓶}}
= -\frac{e^{n\tau_0}}{2n\,(Q^n - Q^{-n})(T^n + T^{-n})} .
\]
したがって AM の「奇数 $n$ の $c = 1$ ＋ Klein 瓶」は、全 $n$ の凍結した状態の項 $e^{n\tau_1/2}/(n(Q^n - Q^{-n}))$ と、$\varepsilon = +1$ の不変な状態の項の和に等しい。
Klein 瓶の項とその分子 $Q^n/T^n - T^n/Q^n$（AM の論文の変数では $q/t - t/q$）は、独立な多重項ではない。奇数の $d$ で整理したときの $-\tfrac12\psi^2 C$ と $\tfrac12 E/(1+t)$ の差である（$E \ne \psi^2 C$）。
\end{proposition}
\begin{proof}
左辺を通分すると、$X^{2n}$ の係数は $[-(T^n + T^{-n}) + Q^{-n}(Q^nT^{-n} - T^nQ^{-n})]/(\cdots) = -T^nQ^{-n}(Q^n + Q^{-n})/(\cdots)$ になり、$(Q^{2n} - Q^{-2n}) = (Q^n - Q^{-n})(Q^n + Q^{-n})$ で約分できる。
\end{proof}

\begin{remark}
AM と AgSc の規約は、$(q, t, Q) \to (q^{-1}, t^{-1}, Q^{-1})$ と全体の符号（$\log Z$ と $-\log Z$）で移り合う。凍結した状態と向き付けられた状態の電荷の比は、AgSc で $x^2/Q_0 = w^{-1}$、AM で $w$ である（$w = (q/t)^{1/2}$）。
閉弦では状態が P$^1$ と RP$^2$ の二つだけで打ち消しが無いので、$E$ と $C$ は一意に決まる。命題 \ref{prop:EC} の不定性は、開弦で初めて問題になる。
\end{remark}

% [O8]
\begin{remark}[取っ手の因子]
AgSc §4 の強い整数性の導出に現れる「差 $\times$ 和」$(q^{1/2} - q^{-1/2})(q^{1/2} + q^{-1/2})$ は、類 $\beta/2$ の $c = 1$ の寄与を $d$ を倍にして並べたときの $\psi^2 z$ である。これは取っ手一つの因子ではない。
Piazzalunga--Uranga \cite{PU14} の M 理論の解釈では、$\tilde I$ で不変な被覆の取っ手一つの捻ったトレースは、同じ変数で $z$ 一つになる。
refined では $G = z_q z_t = T_+T_-$、$T_\pm = (V - V^{-1}) \pm (W - W^{-1})$ と分解できる。ただし $T_\pm$ を使っても、$\sigma$ 対称でない $c = 1$ 部分は展開できるようにならない。
$\mathbb{Z}[W^{\pm 2}][T_+, T_-] = \mathbb{Z}[W^{\pm2}][V - V^{-1}, W - W^{-1}]$ の元のうち $V$ の偶数冪の部分は、$(V - V^{-1})$ の偶数冪の和なので、$V \to V^{-1}$（$\sigma$）で不変だからである。
\end{remark}
